In curvilinear motion of a planet around a star, the direction of the velocity
vector changes continuously. This can be represented by a so-called
``trajectory in velocity space'' and is obtained as follows: for each point on
the spatial trajectory, the corresponding velocity vector is drawn so that its
starting point is at the origin of the velocity space, and its magnitude and
direction is the same as the velocity vector at that point. The tip of this
variable velocity vector generates a curve in velocity space. (The name
`hodograph' was given to this curve by Hamilton in 1846.)

As an example, see figures 1 and 2 below. For a circular orbit (Figure 1), the
magnitude of the velocity is constant and therefore, the hodograph (Figure 2)
of the velocity vector for Keplerian circular motion is also a circle, the
center of which is located at the origin of the velocity space. The radius of
this circle is equal to the constant magnitude of the circular velocity.

\ioaafig{2021_TH_Q12-f1.pdf}

\begin{description}
\item[(12.1)] Write an expression for the radius of the hodograph in Fig.~2,
  as a function of the mass $M$ of the star, and the radius $R$ of the
  circular orbit of the planet's motion. \hfill[1.0\,pt]
\item[(12.2)] For a planet in a Keplerian trajectory, write the expression for
  centripetal acceleration vector ($\vec{a}$) and the magnitude of angular
  momentum ($L$). For any Keplerian trajectory, it is true that
  \[ |\Delta v| = k\Delta\theta \tag{1} \]
  Where $k$ is a constant for each type of Keplerian trajectory. Find the
  expression for the constant $k$ as a function of the masses $M$ and $m$ of
  the star and the planet, respectively, and the angular momentum, $L$.\\
  (Eq.~1) allows us to conclude that for any Keplerian trajectory, the
  hodograph ($v$ as a function of $\theta$) is a circle, but except for
  circular motion, the centre of the hodograph does not coincide with the
  star. It is not necessary to prove this result, you may simply accept it as
  a given. For the hodograph of uniform circular motion, the compliance with
  (eq.~1) is completely obvious, as evidenced in Fig.~3. \hfill[4.0\,pt]

  \ioaafig{2021_TH_Q12-f2.pdf}

\item[(12.3)] Determine the expression of the constant $k$ for the hodograph
  of circular planetary motion. \hfill[2.0\,pt]
\item[(12.4)] Given that the hodograph of the Keplerian elliptical motion is a
  circle, determine the radius of this hodograph and the distance between the
  center of the hodograph and the position of the star, as a function of the
  velocities at periastron and apoastron. Draw a rough sketch of the hodograph
  in the answer sheet as per the schematic shown in Fig.~4. The black circle
  is the star. \hfill[4.0\,pt]

  \ioaafig{2021_TH_Q12-f3.pdf}

\item[(12.5)] Similarly, for the parabolic Keplerian trajectory, determine the
  radius of the corresponding hodograph and the distance from the center of
  that hodograph circle to the star. Express the radius as a function of the
  velocity at periastron. Draw a rough sketch of the hodograph circle in the
  answer sheet. \hfill[4.0\,pt]
\end{description}
